\scanpage{23}
% @ upper left, letterhead
% letterhead
M. DUFOUR\\
VARITYPISTE\\
131, \textsc{Rue Lecourbe}\\
PARIS XV\textsuperscript{e}\\
$\blacksquare$\\
C. C. P. Paris 14.143-34\\
R. M. 140.029

% @ 1–4, top
$Y^{*}_{\deleted{\unsure{1}}}$, Donc maintenant on est dans le cas où la descente de $X'{}^{*}$ donne un $Y'{}^{*}$ complètement décomposé\deleted{\unsure{s}}.

% @ 5–10, above fold
Il en résulte aisément que les composantes connexes de $X'{}^{*}$ sont les images inverses des composantes connexes de $X^{*}$ (\uline{\unsure{noter|cf.} Lemme} 3] compte tenu que \deleted{$\unsure{X\to|X\leftarrow|X}$} $Y'{}^{*}\to Y^{*}$ est univ.\textsuperscript{t} ouvert. Donc

% @ 11–14, bracketed lemma
\uline{Lemme 4} Il existe un rev. étale surjectif $Y^{*}\longrightarrow Y$ tel que $X^{*}=X\times_Y Y^{*}$ ait des composantes connexes qui sont de degré $1$ sur $Y^{*}$.

% @ 15–19, bottom
Le morphisme canonique
\[
X^{*}\longrightarrow Y^{*}\times I=X^{*}_0
\]
[$I$ ens des composantes connexes de $X^{*}$] est compatible \deleted{\unsure{trivialemt}} \inserted{par définition} avec les \unsure{structures|données} de descente sur $Y^{*}\times I$, (et $X^{*}$), \unsure{\textbackslash}
% @ sideways watermark through the text
% printed form
\textsc{Extra Strong}\\
\unsure{Guerimand|Guérimand} Voiron
% archive sheet number: 22
